%% kapa: exam=1 round=32 year=2021 period=1 form=A subject=economics subject_ko=경제학원론 seq=14 number=54
%% source: Q-Net 감정평가사 기출문제 (cst00309) artlSeq=5208397 file=extracted/1008/5208397/A형/2021년도 제32회 감정평가사 1차 1교시 A형.pdf
%% conversion: kordoc 4.24.51 -> tools/gen1.py (auto); review=proofread-sonnet; flags=none
\begin{question}{54}
역선택에 관한 설명으로 옳은 것은?
\begin{choices}
  \item 동일한 조건과 보험료로 구성된 치아보험에 치아건강상태가 좋은 계층이 더 가입하려는 경향이 있다.
  \item 역선택은 정보가 대칭적인 중고차시장에서 자주 발생한다.
  \item 역선택 방지를 위해 통신사는 소비자별로 다른 요금을 부과한다.
  \item 의료보험의 기초공제제도는 대표적인 역선택 방지 수단이다.
  \item 품질표시제도는 역선택을 방지하기 위한 수단이다.
\end{choices}
\end{question}
\answer{5}
